\documentclass[10pt,a4paper]{article}

\usepackage[T1]{fontenc}
\usepackage{amsmath}
\usepackage[margin=2cm]{geometry}
\usepackage{booktabs}
\usepackage{tabularx}
\usepackage{enumitem}
\usepackage{titlesec}
\usepackage[hidelinks]{hyperref}
\usepackage{microtype}

\setlist{nosep,leftmargin=*}
\titlespacing*{\section}{0pt}{1.1ex plus .2ex}{0.5ex}
\titleformat{\section}{\normalfont\large\bfseries}{\thesection.}{0.5em}{}
\renewcommand{\arraystretch}{1.1}
\newcolumntype{L}{>{\raggedright\arraybackslash}X}
\pagestyle{empty}

\begin{document}

\begin{center}
{\large\bfseries Pengaruh Mesin OCR dan Embedding Teks terhadap Pemanfaatan\\
Modalitas Teks pada Klasifikasi Dokumen Multimodal}\\[0.4em]
{\small Proposal Tugas Akhir (revisi) --- Analisis dan Pemrosesan Data Multimodal}\\[0.2em]
{\small Ramadhani Al-Qadri \quad$\cdot$\quad 9 Oktober 2026}
\end{center}
\vspace{0.3em}

\section{Latar belakang dan celah penelitian}

Dokumen hasil pindai membawa informasi di dua lapis: \textbf{tata letak} yang
terlihat pada citra, dan \textbf{isi tekstual} yang hanya bisa diambil lewat OCR.
Karena teksnya diturunkan dari citra, \textbf{mutu cabang teks dibatasi oleh mutu
OCR-nya}.

\textbf{Acuan terkini.} Kashyap \textit{et al.}~\cite{kashyap2026} menggabungkan
EfficientNetB3 (citra) dan BERT (teks OCR) dengan rata-rata logit tanpa bobot, dan
melaporkan akurasi 92,00\% pada Tobacco3482. Ada empat catatan kritis:
\begin{enumerate}[label=(\alph*)]
  \item Mesin OCR-nya hanya disebut ``mesin OCR komersial'' tanpa nama, sehingga
  cabang teksnya tidak bisa direproduksi.
  \item Model dilatih pada 16 kelas RVL-CDIP, padahal Note dan Report (466 dokumen)
  tidak ada di RVL-CDIP. Bila 3.482 dokumen dievaluasi, akurasi maksimum yang mungkin
  hanya 86,6\%. Jadi angka 92\% hampir pasti dihitung pada sebagian kelas saja.
  \item Hasil itu dibandingkan dengan Audebert \textit{et al.}~\cite{audebert2020},
  yang dilatih hanya pada 800 dokumen Tobacco3482, tanpa catatan bahwa protokolnya
  berbeda.
  \item Satu kali \textit{run}, tanpa uji statistik.
\end{enumerate}
Penulisnya sendiri menyebut degradasi OCR yang parah dan pembobotan modalitas adaptif
sebagai arah lanjutan.

\textbf{Gambaran lapangan.} Meta-analisis atas 139 studi~\cite{mironczuk2026}
menemukan fusion multimodal menaikkan akurasi rata-rata 5,28 poin. Namun hanya 11,8\%
studi multimodal yang memakai uji statistik, dan kualitas OCR \textit{tidak dapat
dikontrol} dalam meta-analisis itu karena studi primer jarang melaporkannya.
Perbandingan mesin OCR~\cite{francis2025} mengukur ketepatan OCR itu sendiri, bukan
dampaknya pada tugas lanjutan. Padahal galat OCR terbukti merambat ke tugas
hilir~\cite{zhang2025}.

\textbf{Hasil awal proyek ini.} Replikasi Audebert pada Tobacco3482 (Tesseract +
FastText, 3 \textit{split} acak, 800 dokumen latih) mencapai angka berikut:

\begingroup\small
\begin{center}
\begin{tabular}{@{}lcc@{}}
\toprule
Model & Akurasi kami & Akurasi paper \\
\midrule
Teks (FastText + CNN1D) & $0{,}727 \pm 0{,}006$ & 0,738 \\
Citra (MobileNetV2) & $0{,}809 \pm 0{,}021$ & 0,845 \\
Fusion (konkatenasi) & $0{,}834 \pm 0{,}011$ & 0,878 \\
Oracle (teks \textbf{atau} citra benar) & 0,898 & 0,921 \\
\bottomrule
\end{tabular}
\end{center}
\endgroup

\noindent Fusion tampak unggul. Tetapi \textbf{menolkan seluruh masukan teks hanya
menurunkan akurasi fusion sebesar 0,0015}: cabang teksnya praktis tidak terpakai.
Diagnostik kami menelusuri penyebabnya sampai ke mekanisme: fitur teks kolaps (67 dari
128 unit mati) dan skalanya 3,7$\times$ lebih kecil daripada fitur citra. Setelah
diperbaiki (LayerNorm per cabang dan inisialisasi dari model teks; seed 42), menolkan
teks menurunkan akurasi 0,689 (sebelumnya 0,007), dan F1 kelas Resume naik dari 0,81 ke
0,97. Akurasinya turun 0,019, tetapi macro F1 tetap identik. Pelajarannya:
\textbf{akurasi fusion yang tinggi bukan bukti bahwa fusion memakai kedua modalitas.}

\medskip
\noindent\fbox{\parbox{\dimexpr\textwidth-2\fboxsep-2\fboxrule}{\textbf{Celah
penelitian.} Studi fusion teks--citra memperlakukan mesin OCR sebagai konstanta,
bahkan tanpa menyebut namanya, dan hanya melaporkan akurasi gabungan. Belum diukur
bagaimana \textbf{pilihan mesin OCR $\times$ jenis embedding} memengaruhi kekuatan
cabang teks, dan \textbf{seberapa jauh model fusion benar-benar memakai cabang
itu}.}}

\section{Pertanyaan penelitian}

\begin{enumerate}[label=\textbf{Q\arabic*.},leftmargin=2.4em]
  \item Seberapa besar pilihan mesin OCR (Tesseract, EasyOCR, PaddleOCR, Keras-OCR)
  mengubah akurasi cabang teks pada Tobacco3482?
  \item Apakah embedding kontekstual (BERT) memperkecil selisih antar-mesin OCR
  dibanding embedding subword statis (FastText)? Dengan kata lain, adakah interaksi
  OCR $\times$ embedding?
  \item Apakah cabang teks yang lebih kuat membuat model fusion lebih \textit{memakai}
  teks, diukur dengan uji \textit{missing modality}? Dan apakah aturan rata-rata logit
  Kashyap lebih tahan terhadap cabang yang timpang daripada fusion fitur?
\end{enumerate}

\section{Desain eksperimen}

\begingroup\small
\begin{tabularx}{\textwidth}{@{}p{2.7cm} L@{}}
\toprule
\textbf{Faktor} & \textbf{Level} \\
\midrule
Mesin OCR & Tesseract~4 (\texttt{-{}-oem 1 -{}-psm 3}), EasyOCR, PaddleOCR; Keras-OCR
bila dapat dijalankan. \textbf{Semua dijalankan pada citra JPG yang sama.}
QS-OCR (Tesseract pada TIF asli) dipakai sebagai acuan untuk mengukur efek JPG vs TIF. \\
Embedding teks & FastText 300-d + CNN~1D~\cite{bojanowski2017} (pipeline hasil awal);
BERT-base~\cite{devlin2019} di-\textit{fine-tune}, seperti Kashyap. Ketahanan
embedding terhadap derau OCR masih merupakan persoalan terbuka~\cite{michail2025}. \\
Cabang citra & MobileNetV2~\cite{sandler2018}, $384\times384$, \textbf{tetap} di semua
kondisi, supaya perbedaan hanya berasal dari teks. \\
Fusion & (i) rata-rata logit tanpa bobot, aturan Kashyap, tanpa pelatihan ulang;
(ii) fusion fitur dengan LayerNorm per cabang, hasil perbaikan pada tahap awal. \\
\bottomrule
\end{tabularx}
\endgroup

\smallskip
Teks masuk ke kedua embedding \textbf{tanpa} praproses gaya Kashyap (membuang angka
dan \textit{stopword}), agar efek embedding tidak tercampur efek praproses. Selain
teks, setiap mesin OCR menghasilkan \textbf{kotak posisi teks}. Kotak ini disimpan
sejak awal untuk tahap lanjutan (Bagian~6), sehingga OCR tidak perlu dijalankan ulang.

\section{Dataset dan lisensi}

\textbf{Tobacco3482-jpg} (3.482 citra, 10 kelas; Kaggle
\texttt{patrickaudriaz/tobacco3482jpg}) tidak memiliki lisensi baku; deskripsinya
hanya merujuk Kumar \textit{et al.}~\cite{kumar2014}. \textbf{QS-OCR-Small}
(GitHub \texttt{QuickSign/ocrized-text-dataset}, rilis \texttt{v1.0}) berlisensi
Apache-2.0, tetapi teksnya turunan Tobacco3482 sehingga hak sumber hulu tetap
berlaku. Pemakaian dibatasi pada konteks akademik non-komersial. Tidak ada berkas
dataset yang diredistribusikan: notebook mengunduh langsung dari sumber aslinya.
RVL-CDIP tidak dipakai karena ukurannya ($\sim$400.000 dokumen) di luar anggaran
komputasi. Konsekuensinya, angka kami \textbf{tidak sebanding langsung} dengan
Kashyap, yang dilatih pada RVL-CDIP.

Protokolnya mengikuti hasil awal: 800 dokumen latih per \textit{split},
terstratifikasi, dengan seed 42/43/44, dan indeks \textit{split} yang sama untuk
semua kondisi.

\section{Metrik evaluasi}

\begingroup\small
\begin{tabularx}{\textwidth}{@{}p{3.9cm} L@{}}
\toprule
\textbf{Metrik} & \textbf{Alasan} \\
\midrule
Akurasi (OA), macro F1, F1 per kelas & OA untuk perbandingan dengan literatur. Macro F1
karena kelasnya timpang (5{,}2:1). Pada hasil awal, OA dan macro F1 terbukti bisa
bergerak terpisah. \\
Proksi kualitas OCR & Tobacco3482 tidak punya transkripsi acuan, jadi \textit{character
error rate} tidak dapat dihitung. Yang diukur: jumlah token, laju kata di luar
kosakata (OOV), dan kesepakatan antar-mesin. Dilaporkan sebagai proksi, bukan akurasi
OCR. \\
Uji \textit{missing modality} & Penurunan akurasi saat teks atau citra dinolkan.
Inilah metrik yang menjawab Q3. \\
Oracle & Batas atas teoretis: benar bila salah satu cabang benar. \\
Uji statistik & McNemar berpasangan antar-kondisi pada test set yang sama, ditambah
rata-rata $\pm$ simpangan baku 3 seed. Menjawab kritik metodologis~\cite{mironczuk2026}. \\
\bottomrule
\end{tabularx}
\endgroup

\section{Tahap lanjutan: teks sebagai objek}

Kotak teks dari Bagian~3 memungkinkan teks diperlakukan sebagai \textbf{objek
berposisi}, bukan hanya rangkaian kata. Representasinya bisa berupa fitur posisi 2D di
samping embedding kata, peta kotak teks sebagai citra, atau graf antar-kotak. Satu
kontrol bersifat wajib: dokumen Tobacco3482 bercap \textbf{kode ID} arsip yang
posisinya konsisten, dan fitur dari kode ID saja sudah memberi $\pm$60\% akurasi di
Tobacco3482~\cite{larson2025}. Model berbasis posisi harus diuji dengan dan tanpa kotak
kode ID. Tanpa kontrol itu, akurasi tinggi tidak dapat ditafsirkan.

\section{Kelayakan dan risiko}

Komputasi memakai laptop RTX 3050 4~GB VRAM. OCR diperkirakan memakan $\pm$1--3 jam per
mesin untuk 3.482 halaman. BERT-base dilatih dengan \textit{mixed precision} dan
akumulasi gradien; bila tidak muat, diganti DistilBERT dan dicatat sebagai deviasi.
Rata-rata logit tidak membutuhkan pelatihan, sehingga dapat diuji untuk semua kondisi.
Fusion fitur hanya dilatih untuk kondisi OCR terbaik dan terburuk.

Risiko teknisnya:
\begin{enumerate}[label=(\alph*)]
  \item Tesseract di Windows bukan paket pip.
  \item \texttt{paddlepaddle-gpu} berukuran $>$700~MB.
  \item Keras-OCR berbasis TensorFlow, yang sejak versi 2.11 tidak mendukung GPU di
  Windows native, dan rilis terakhirnya November 2023.
\end{enumerate}
Bila Keras-OCR tidak dapat dijalankan, ia dilewati dengan alasan tertulis. Tiga mesin
lainnya tetap menjawab Q1--Q3.

\begingroup\footnotesize
\begin{thebibliography}{99}
\bibitem{kashyap2026} A. Kashyap, M. Mittal, S. Singh, ``Multimodal fusion and
supervised contrastive learning for robust document image classification,''
\textit{International Journal of Intelligent Systems}, vol.~2026, art.~4241437, 2026.
doi:10.1155/int/4241437.
\bibitem{mironczuk2026} M. M. Miro\'nczuk, ``Document classification pattern
recognition via information fusion: A systematic review of multimodal and multiview
representation approaches,'' \textit{Information Fusion}, vol.~132, art.~104247, 2026.
doi:10.1016/j.inffus.2026.104247.
\bibitem{francis2025} S. A. Francis, M. Sangeetha, ``A comparison study on optical
character recognition models in mathematical equations and in any language,''
\textit{Results in Control and Optimization}, vol.~18, art.~100532, 2025.
doi:10.1016/j.rico.2025.100532.
\bibitem{zhang2025} J. Zhang \textit{et al.}, ``OCR hinders RAG: Evaluating the
cascading impact of OCR on retrieval-augmented generation,'' in \textit{Proc.\ IEEE/CVF
ICCV}, 2025, pp.~17443--17453. doi:10.1109/ICCV51701.2025.01620.
\bibitem{michail2025} A. Michail, J. Opitz, Y. Wang, R. Meister, R. Sennrich,
S. Clematide, ``Cheap character noise for OCR-robust multilingual embeddings,'' in
\textit{Findings of ACL}, 2025.
\bibitem{larson2025} S. Larson, S. Duwal, B. Vilnrotter, G. Chakkithara, V. Padwal,
K. Leach, ``Spurious cues in RVL-CDIP and Tobacco3482 document classification: The case
of ID codes,'' in \textit{Proc.\ ACM DocEng}, 2025. doi:10.1145/3704268.3748683.
\bibitem{audebert2020} N. Audebert, C. Herold, K. Slimani, C. Vidal, ``Multimodal deep
networks for text and image-based document classification,'' in \textit{ECML PKDD 2019
Workshops}, CCIS vol.~1167, Springer, 2020, pp.~427--443.
doi:10.1007/978-3-030-43823-4\_35.
\bibitem{kumar2014} J. Kumar, P. Ye, D. Doermann, ``Structural similarity for document
image classification and retrieval,'' \textit{Pattern Recognition Letters}, vol.~43,
2014.
\bibitem{devlin2019} J. Devlin, M.-W. Chang, K. Lee, K. Toutanova, ``BERT: Pre-training
of deep bidirectional transformers for language understanding,'' in \textit{Proc.\
NAACL-HLT}, 2019.
\bibitem{bojanowski2017} P. Bojanowski, E. Grave, A. Joulin, T. Mikolov, ``Enriching
word vectors with subword information,'' \textit{TACL}, vol.~5, 2017.
\bibitem{sandler2018} M. Sandler \textit{et al.}, ``MobileNetV2: Inverted residuals and
linear bottlenecks,'' in \textit{Proc.\ CVPR}, 2018.
\end{thebibliography}
\endgroup

\end{document}
